Three-dimensional (3D) assets are central to games, augmented and virtual reality, and online retail, yet producing them by hand is slow and expensive. Recent image-to-3D methods can create a mesh from a single photograph, but they produce dense meshes with tens to hundreds of thousands of faces, which are too heavy for real-time use. Standard mesh simplification such as Quadric Error Metrics (QEM)~\cite{garland1997qem} reduces the face count but ignores the texture, so colours stretch and tear along UV seams.

GEN3D addresses both problems in one pipeline. A generator that we design and train ourselves turns one photograph into a textured 3D mesh. It combines a pretrained DINOv2 encoder~\cite{oquab2023dinov2} with pixel-aligned features~\cite{saito2019pifu,xu2019disn}, an occupancy decoder~\cite{mescheder2019onet} and a texture field~\cite{oechsle2019texturefields}. The resulting high-poly mesh is then simplified by a texture-aware QEM and, as the main contribution of the project, by D-QEM, in which a Transformer~\cite{vaswani2017attention} learns a correction to the QEM collapse cost that accounts for UV seams, UV stretch and colour change.

This report covers the first 40\% of the project. We built a class-agnostic dataset of 7{,}467 clean objects from 1{,}028 categories of Objaverse-LVIS~\cite{deitke2023objaverse}, wrote our own inside/outside labelling method for non-watertight meshes, rendered 16 calibrated views per object with Blender, automated all heavy computation on Kaggle GPUs, and implemented and verified the generator. The texture-aware QEM, the D-QEM Transformer and the interactive demo follow in the remaining modules.
